% This file is intended to replace the main .tex file in the official ACL Overleaf template.
% Also replace the template's custom.bib with the accompanying custom.bib file.
% Keep acl.sty and acl_natbib.bst from the official template in the project root.

\documentclass[11pt]{article}

% Overleaf speed switches.
% Fast editing mode: set \draftfigurestrue and \includeappendixfiguresfalse.
% Final PDF mode: keep \draftfiguresfalse and \includeappendixfigurestrue.
\newif\ifdraftfigures
\draftfiguresfalse
\newif\ifincludeappendixfigures
\includeappendixfigurestrue

% Change "review" to "final" for camera-ready.
\usepackage[review]{acl}

% Standard package includes.
\usepackage{times}
\usepackage{latexsym}
\usepackage[T1]{fontenc}
\usepackage[utf8]{inputenc}
\usepackage{microtype}
\usepackage{inconsolata}

% Additional packages for figures and tables.
\ifdraftfigures
  \usepackage[draft]{graphicx}
\else
  \usepackage{graphicx}
\fi
\usepackage{booktabs}
\usepackage{tabularx}
\usepackage{multirow}
\usepackage{adjustbox}
\usepackage{amsmath}
\usepackage{array}
\usepackage{xurl}
\usepackage{placeins}

\newcommand{\ToolRobustBench}{\textsc{ToolRobustBench}}
\newcommand{\safeincludegraphics}[2][]{%
  \IfFileExists{#2}{\includegraphics[#1]{#2}}{%
    \IfFileExists{../#2}{\includegraphics[#1]{../#2}}{%
      \fbox{\parbox{0.95\linewidth}{\centering Missing figure file: \\ \texttt{\detokenize{#2}}}}%
    }%
  }%
}
\newcolumntype{Y}{>{\raggedright\arraybackslash}X}

\title{\ToolRobustBench: Stage-Wise Perturbation Evaluation and Failure Diagnosis for Tool-Calling Robustness}

\author{Anonymous ACL submission}

\begin{document}
\maketitle

\begin{abstract}
Large language models increasingly rely on function calling and tool execution, yet aggregate success rates under clean conditions do not explain where tool-use failures begin or how they cascade across a call. We introduce \ToolRobustBench, a stage-wise diagnostic benchmark for tool-calling robustness. It aligns four perturbation families with the tool-use pipeline: tool-interface perturbation, user-intent perturbation, tool-output / observation perturbation, and runtime-environment perturbation, and attributes failures across tool selection, schema grounding, argument binding, runtime-feedback handling, and end-to-end task success.

We evaluated 15,456 single-family instances across 7 models, 16 sampled local tools, 4 perturbation families, and 14 subtypes in a deterministic local-tool environment. Clean-setting performance is near saturated, with an average success rate of 97.9\%, but family-level robustness drops to 70.8\% under perturbation and 49.1\% under heavy perturbation. The main bottleneck is tool-output / observation perturbation, where no model reaches 0.60, and the average success rate is 45.5\%. Representative cross-family experiments further show that composition difficulty is not reducible to the weakest single-family component: A3+C4 and B4+C4 create strong negative interactions, while A2+B3 and B2+D2 expose floor effects or apparent reversals. \ToolRobustBench provides a reproducible benchmark for measuring robustness, diagnosing cascade-aware error sources, and improving tool-using agents beyond clean tool-call accuracy.
\end{abstract}

\section{Introduction}

Tool calling has become a central interface through which large language models move from text generation to executable agents. Unlike ordinary question answering, tool use requires a multi-stage decision process: the model must first understand the user's intent, then choose the correct tool from a candidate set, bind the request to structured arguments, and finally make a stable judgment from tool returns or execution feedback. In real deployments, this process is often noisy. Tool names may change from semantic labels to semi-opaque service names; descriptions of nearby tools may overlap heavily; user requests may be indirect, underspecified, or internally conflicting; and tool feedback may contain schema drift, evidence loss, false facts, misleading candidates, or runtime failures. Therefore, tool-use evaluation should not only measure final accuracy under clean inputs, but also test whether model behavior remains robust, interpretable, and diagnosable under perturbed environments.

Existing work has advanced tool-use evaluation from several directions. ToolBench, API-Bank, and StableToolBench study large-scale tool use, API calling ability, and reproducible execution environments, providing important foundations for evaluating tool selection and call formatting \citep{qin2024toolllm,li2023apibank,guo2024stabletoolbench}. ToolSandbox and $\tau$-bench extend evaluation from single calls to process-level and trajectory-level interaction, emphasizing state tracking and behavioral diagnosis in multi-step environments \citep{lu2025toolsandbox,yao2025taubench}. AgentNoiseBench injects realistic user-side and tool-side noise while preserving task solvability \citep{wang2026agentnoisebench}. Together, these studies show that tool-use robustness is a core problem for agent evaluation, and that noise injection and process diagnosis provide a finer-grained view than clean accuracy alone.

However, a key diagnostic gap remains. When a model fails under perturbed tool use, we need to know not only whether it failed, but why the failure occurred, which stage failed first, what the primary failure type was, and whether the final observed error was merely a cascade from an upstream mistake. This distinction is crucial for tool calling. A final wrong-parameter failure does not necessarily mean that argument binding was directly perturbed; it may originate from earlier tool-selection confusion or from an interface description that pushed the model toward the wrong tool path. Likewise, an execution failure may not be the direct result of runtime perturbation, but a downstream consequence of wrong arguments, wrong tools, or wrong interpretation. If a benchmark records only the final failure label, it can mistake reasonable cascade behavior for perturbation-construction error and can obscure the stage where the model is truly fragile.

To address this gap, we propose \ToolRobustBench, a controlled robustness benchmark for single-step tool calling. Unlike tool-use benchmarks that mainly ask whether the final task succeeds, \ToolRobustBench{} decomposes the tool-calling process into four perturbation families: tool-interface exposure, user intent, tool-output / observation return, and runtime environment. Each family contains calibrated light, medium, and heavy perturbation severities. More importantly, \ToolRobustBench{} explicitly tracks observed failure, primary failure, first failure stage, cascade source, and boundary violation. This lets us determine whether a failure is the direct effect of the intended perturbation family or a natural cascade from an upstream stage. The evaluation goal is therefore expanded from ``does the model get the noisy case right?'' to ``where does the model first go wrong, how does the error propagate, and whether the perturbation boundary remains clear.''

This paper addresses that gap with \ToolRobustBench, a reproducible and attributable benchmark for robust tool calling. Our main contributions are:

\begin{enumerate}
    \item \textbf{Stage-aligned perturbation taxonomy.} We define four perturbation families: tool-interface, user-intent, tool-output / observation, and runtime-environment perturbation, that correspond to distinct entry points in the tool-calling pipeline.
    \item \textbf{Cascade-aware error attribution.} We distinguish observed error labels from primary error labels, earliest failure stages, error source stages, and boundary-violation indicators, so downstream symptoms are not mistaken for direct perturbation failures.
    \item \textbf{Reproducible deterministic local-tool benchmark.} We evaluate sampled local tools from a 40-tool deterministic environment, enabling stable gold arguments, expected results, and failure attribution without live-API volatility.
    \item \textbf{Empirical findings on single-family and cross-family robustness.} Across 15,456 single-family instances and representative mixed perturbations, we show that observation handling is the main bottleneck and that cross-family difficulty is not reducible to the weakest single-family component.
\end{enumerate}

Overall, \ToolRobustBench{} complements existing tool-use benchmarks by making robustness failures localizable, reproducible, and interpretable across stages.

\begin{figure*}[t]
    \centering
    \safeincludegraphics[width=0.95\linewidth]{picture/fig01_toolrobustbench_pipeline.png}
    \caption{\ToolRobustBench{} diagnostic pipeline. The figure shows the relationship among the user request, tool registry, model tool call, tool execution, runtime feedback, post-execution decision, and scorer, and marks the insertion positions of the four perturbation families A/B/C/D.}
    \label{fig:pipeline}
\end{figure*}

\section{Related Work}

\paragraph{Tool-calling evaluation.}
API-Bank and ToolBench extend evaluation from question answering to API retrieval, tool selection, and argument calling \citep{li2023apibank,qin2024toolllm}. StableToolBench shows that live API instability can contaminate evaluation and therefore uses stable simulated APIs \citep{guo2024stabletoolbench}. \ToolRobustBench{} follows the reproducibility motivation but evaluates controlled perturbations across tool interfaces, user intent, tool outputs, and runtime environments rather than clean API calling alone.

\paragraph{Agent-process evaluation.}
ToolSandbox and $\tau$-bench focus on stateful and interactive tool use, where evaluation depends on state updates, action sequences, and trajectory quality \citep{lu2025toolsandbox,yao2025taubench}. \ToolRobustBench{} remains single-step, but the tool-output and runtime-environment families introduce a post-execution decision stage: the model first makes a tool call, receives perturbed feedback, and must recover the result or decide the execution state. This isolates observation interpretation and runtime-state judgment without requiring a long-horizon environment.

\paragraph{Noise-robustness evaluation.}
AgentNoiseBench categorizes realistic user-side and tool-side noise with solvability-preserving perturbations \citep{wang2026agentnoisebench}, while RoTBench studies multi-level robustness for tool learning \citep{ye2024rotbench}. Our difference is diagnostic granularity: \ToolRobustBench{} binds perturbation families to tool-use stages and reports cascade-aware attribution fields, allowing each failed instance to be traced through an observed error label, primary error label, earliest failure stage, and boundary-violation indicator.

\section{Design of \ToolRobustBench}

\subsection{Tool-Calling Capability Axes}

\ToolRobustBench{} decomposes tool use into five diagnostic capability axes. They follow the temporal order of a tool call, but their main role is attribution: each axis identifies a capability that can become the earliest source of failure.

\begin{table*}[t]
\centering
\small
\begin{tabularx}{\linewidth}{l l Y Y}
\toprule
Stage & Name & Target capability & Typical failures \\
\midrule
S1 & Tool Selection & Select the correct tool from candidate tools & no-call, tool misselection \\
S2 & Schema Grounding & Understand tool names, descriptions, and schemas & tool misselection, schema-format violation \\
S3 & Argument Binding & Bind user intent to arguments & missing argument, incorrect argument \\
S4 & Tool Output / Runtime Feedback Handling & Interpret tool returns and runtime feedback & incorrect argument, schema-format violation, runtime failure \\
S5 & End-to-End Task Success & Whether the end-to-end task succeeds & all failure tags \\
\bottomrule
\end{tabularx}
\caption{Five diagnostic capability axes in \ToolRobustBench.}
\label{tab:stages}
\end{table*}

This decomposition provides the coordinate system for perturbation design and failure attribution. Because downstream behavior may be a cascade, each record preserves the earliest failure stage, primary error label, observed error label, and error source stage.

\subsection{Stage-Aligned Perturbation Taxonomy}

\ToolRobustBench{} contains four single-family perturbation types aligned with these axes.

\textbf{Tool-interface perturbation} simulates interface-level shifts, including surface changes in tool names, abstract descriptions, and blurred boundaries among similar tools. This family tests whether models over-rely on name or description cues.

\textbf{User-intent perturbation} simulates user-expression shifts, including indirectness, underspecification, distracting context, and internal conflicts. This family tests whether models can recover task goals and required arguments from natural language.

\textbf{Tool-output / observation perturbation} simulates problems in processing perturbed returned content. After the initial call, the model must recover the clean target from noisy feedback, distinguish the primary result from supporting evidence, and handle missing or conflicting evidence.

\textbf{Runtime-environment perturbation} simulates execution-state and runtime-feedback shifts. The model must decide whether the output is success, partial success, or error, and recover the corresponding retryability, failure type, and final result.

At the template level, the single-family setting contains 14 subtypes with clean, light, medium, and heavy severities. The pairwise setting contains 5 representative family pairs and 122 cross-family templates at medium or heavy severity; full template details are deferred to the appendix.

\subsection{Cascade-Aware Error Attribution}

Real tool-calling errors often have a cascade structure. A final incorrect-argument symptom may originate from S1/S2 tool misselection; a runtime failure may be induced by upstream argument corruption rather than the execution environment itself. To avoid confusing downstream symptoms with earliest causes, \ToolRobustBench{} records five diagnostic fields during scoring.

\begin{table*}[t]
\centering
\small
\begin{tabularx}{\linewidth}{l Y Y}
\toprule
Field & Meaning & Purpose \\
\midrule
Observed error label & Final observed failure type & Describes end-to-end output behavior \\
Primary error label & Main failure type after cascade backtracking & Points to the main and earliest error source \\
Earliest failure stage & First failed stage & Localizes the error to S1--S4 \\
Error source stage & Source stage if the failure is cascaded & Distinguishes upstream causes from downstream symptoms \\
Boundary-violation indicator & Whether the failure violates the expected family boundary & Checks perturbation construction and failure attribution \\
\bottomrule
\end{tabularx}
\caption{Cascade-aware diagnostic fields.}
\label{tab:diagnostic-fields}
\end{table*}

The main text focuses on primary failures, cascade rates, and boundary-violation indicators. This separation lets us report both the final symptom and whether that symptom is consistent with the expected perturbation boundary.

\subsection{Metrics}

We report clean-setting performance, family-level robustness, overall robustness, and robustness drop. Clean-setting performance is end-to-end success under the clean condition. Family-level robustness is the macro-averaged success rate of one perturbation family under light, medium, and heavy conditions. Overall robustness is the macro-average over the four families. Robustness drop compares paired clean and perturbed success. Failure analysis uses primary error labels, earliest failure stages, and boundary-violation indicators; lower-level scorer outputs are reported in the appendix.

\section{Experimental Setup}

\subsection{Data and Models}

\ToolRobustBench{} uses deterministic local tools so that tool outputs, gold arguments, expected results, and failure attribution are stable across runs. The complete environment contains 40 tools organized into 12 similar tool groups; the main experiments sample 16 local tools with a fixed seed to control cost and preserve cross-model comparability. The full tool list is reported in Appendix Table~\ref{tab:tool-env}.

We evaluate 7 models on 4 perturbation families and 14 subtypes. For each model, tool-interface and runtime-environment perturbations contain 480 records each, while tool-output / observation and user-intent perturbations contain 624 records each. The single-family analysis therefore covers 15,456 raw instances. Each family is evaluated under clean, light, medium, and heavy conditions, and all models use the same task set, sampled tool registry, and scorer.

\subsection{Scoring Workflow}

For each instance, the model first generates a tool call. For tool-interface, user-intent, and clean instances, the scorer parses the call and checks tool identity, schema validity, arguments, and execution result. For tool-output and runtime-environment instances, the model first completes the initial call; the system then injects perturbed feedback; the model returns an observation recovery or post-execution decision; and the scorer checks the expected result, status, retryability, and failure type.

Provider-specific parameters are used only for protocol compatibility and do not change task content or scoring criteria. The main findings are robust to scorer normalization and protocol-compatibility adjustments; details are deferred to Appendix~\ref{app:validity}.

\section{Single-Family Results}

\subsection{Overall Robustness and Family Ranking}

The single-family experiment asks whether clean tool-calling performance predicts robustness under controlled perturbations. Table~\ref{tab:leaderboard} shows that it does not. The average clean-setting success rate is 0.979, and three models reach 1.000, yet overall robustness ranges from 0.766 to 0.664 once perturbations are introduced. The main takeaway is therefore not a narrow model ranking, but the gap between clean capability and stage-wise robustness.

Family-level results expose the mechanism behind this gap. Averaged across models, tool-interface perturbation is easiest at 0.918, followed by user-intent perturbation at 0.773 and runtime-environment perturbation at 0.688. Tool-output / observation perturbation is the bottleneck at 0.455, and no model reaches 0.60 in this family. This indicates that current models are better at selecting and grounding tools than at recovering reliable results from corrupted or incomplete returned evidence.

The main exception illustrates why family decomposition is necessary. \texttt{deepseek-v4-flash} reaches 0.958 on tool-interface perturbation and 0.783 on user-intent perturbation, but drops to 0.319 on tool-output / observation perturbation. Strong tool selection therefore does not imply strong observation recovery. Conversely, \texttt{qwen-plus} is competitive on runtime-environment perturbation at 0.713 but weak on observation at 0.380, showing that runtime-state judgment and evidence recovery are separable capabilities.

\begin{table*}[t]
\centering
\small
\begin{adjustbox}{width=\linewidth}
\begin{tabular}{lrrrrrrrl}
\toprule
Model & Clean & Overall & A Interface & B Intent & C Output & D Runtime & Heavy Macro & Main failure \\
\midrule
\texttt{claude-sonnet-4-6} & 1.000 & 0.766 & 0.968 & 0.785 & 0.568 & 0.745 & 0.583 & incorrect argument \\
\texttt{deepseek-v4-pro} & 1.000 & 0.749 & 0.926 & 0.809 & 0.512 & 0.748 & 0.516 & incorrect argument \\
\texttt{gpt-5.4-mini} & 0.984 & 0.712 & 0.856 & 0.781 & 0.503 & 0.708 & 0.469 & incorrect argument \\
\texttt{qwen-plus} & 1.000 & 0.706 & 0.931 & 0.802 & 0.380 & 0.713 & 0.507 & incorrect argument \\
\texttt{gemini-2.5-pro} & 0.979 & 0.692 & 0.900 & 0.748 & 0.481 & 0.637 & 0.454 & incorrect argument \\
\texttt{deepseek-v4-flash} & 0.927 & 0.669 & 0.958 & 0.783 & 0.319 & 0.616 & 0.475 & incorrect argument \\
\texttt{gemini-2.5-flash} & 0.964 & 0.664 & 0.889 & 0.700 & 0.420 & 0.648 & 0.436 & incorrect argument \\
\bottomrule
\end{tabular}
\end{adjustbox}
\caption{Main leaderboard recomputed from raw records. The maximum-drop family for every model is tool-output / observation perturbation.}
\label{tab:leaderboard}
\end{table*}

\begin{figure}[t]
    \centering
    \safeincludegraphics[width=\linewidth]{picture/fig03_family_model_heatmap.png}
    \caption{Family-level robustness heatmap. Tool-output / observation perturbation is the largest-drop family for all models, and no model reaches 0.60 in family C.}
    \label{fig:family-heatmap}
\end{figure}

\subsection{Subtype and Severity Effects}

This experiment asks whether family-level rankings are driven by broad perturbation categories or by a few subtype-specific stressors. The severity trend is clear at the aggregate level: average success falls from 0.979 under clean conditions to 0.869, 0.724, and 0.491 under light, medium, and heavy perturbations. Heavy perturbations therefore expose robustness gaps that are mostly hidden by clean performance.

Subtype results further explain the family-level differences. The hardest subtype is C2 Return Evidence Loss, with a success rate of only 0.142, followed by C1 Return Structure Noise at 0.460 and D2 Runtime Variability at 0.469. By contrast, A2 Description Abstraction and A1 Name Surface Shift reach 0.927 and 0.930, indicating that surface interface changes are usually easier to recover from than missing return evidence.

Within B, B3 Distraction has a success rate of 0.658, clearly lower than B1 Indirectness at 0.801, B4 Conflict at 0.812, and B2 Underspecification at 0.820. This result suggests that distracting information is more likely than simple omission or indirectness to push the model toward a wrong tool or wrong task interpretation. Within D, D2 Runtime Variability has a success rate of 0.469, much lower than D1 Availability Failure at 0.736 and D3 Failure Feedback Quality at 0.858, showing that models have the most difficulty stably judging final state when stale/full, fresh/partial, and multiple attempts are mixed.

Subtype-level non-monotonicity also matters. Static template severity is monotonic, but model accuracy does not always decrease strictly from light to medium to heavy. Some heavy runtime-availability failures, such as timeout or permission denied, are more explicit than ambiguous stale-cache plus partial-fragment states and can be easier to classify. We therefore interpret severity as a controlled stress axis rather than as a guarantee of monotonic model accuracy for every subtype.

\subsection{Failure Taxonomy and Cascade Attribution}

Across all perturbed failures, incorrect argument is the most frequent primary error label, with 2,015 cases. It is followed by tool misselection with 1,051 cases, runtime failure with 921 cases, missing argument with 149 cases, no-call with 121 cases, and schema-format violation with 49 cases. This distribution shows that the main problem in tool-calling robustness is not that models completely fail to call tools. Rather, after entering the tool-calling chain, models convert wrong evidence, wrong tool paths, or wrong runtime interpretations into wrong arguments and wrong final states.

When decomposed by family, failure patterns are highly concentrated: 89.1\% of tool-interface failures are tool misselection, 89.9\% of user-intent failures are tool misselection, 91.7\% of tool-output / observation failures are incorrect arguments, and 97.6\% of runtime-environment failures are runtime failures. This concentration supports the family design in \ToolRobustBench: different perturbation sources induce different error types instead of randomly mixing into the same failure pool. Without failure taxonomy, A/B tool-misselection failures, C incorrect-argument failures, and D runtime failures would be collapsed into an uninterpretable aggregate failure rate.

\begin{figure}[t]
    \centering
    \safeincludegraphics[width=\linewidth]{picture/fig04_failure_taxonomy.png}
    \caption{Primary failure taxonomy by perturbation family. Different families have clearly different failure compositions, showing that aggregate success rate must be interpreted together with failure attribution.}
    \label{fig:failure-taxonomy}
\end{figure}

Cascade-aware attribution explains why these labels should not be read mechanically. Runtime-environment perturbation has the highest cascade rate, 0.638, because runtime feedback lies late in the tool chain and absorbs upstream selection or argument errors. Tool-interface perturbation has a cascade rate of 0.352, often because tool-boundary confusion causes downstream argument or execution symptoms. Tool-output / observation perturbation has a cascade rate of only 0.003, showing that its failures usually occur directly during evidence selection or result recovery. Boundary statistics in Appendix Table~\ref{tab:boundary-stats} further show that a high cascade rate is not the same as perturbation-construction failure: runtime-environment perturbation has a boundary-violation rate of only 0.006.

\subsection{Post-Execution Diagnostics}

Runtime-environment instances require a post-execution decision over status, retryability, failure type, and final result. Averaged across models, retryability matching is highest at 0.949, followed by failure-type matching at 0.917, status matching at 0.783, and result matching at 0.738. Models therefore understand many explicit runtime labels but struggle more with recovering the final result from mixed feedback. Detailed post-execution profiles are deferred to Appendix Figures~\ref{fig:app-stage2} and \ref{fig:app-gemini-profile}.

\section{Mixed Perturbation Results}

This section asks whether multi-source perturbations can be predicted from their single-family components. We evaluate five representative pairwise combinations under medium and heavy severities: A2+B3, A3+C4, A1+D1, B4+C4, and B2+D2. Each model, combination, and severity contains 16 perturbed instances, producing 70 comparable cells. Figure~\ref{fig:cross-vs-single} compares cross-family success with the weaker of the two corresponding single-family components.

Overall, cross-family success rates are clearly lower than single-family success rates. Under medium severity, the average success rate over the five combinations and seven models is 0.448, while the average weakest single-family component is 0.621 and the mean of the two components is 0.754. Under heavy severity, the cross-family average further drops to 0.200, while the weakest-component average is 0.292 and the component mean is 0.507. In other words, cross-family difficulty is not simply the minimum of two single-family scores; it often produces additional composition loss. Relative to the mean of the two components, the average gaps under medium and heavy are -0.306 and -0.307.

\begin{figure}[t]
    \centering
    \safeincludegraphics[width=\linewidth]{picture/fig06_crossfamily_vs_single_family_delta.png}
    \caption{Cross-family perturbations versus the weakest single-family component. Gray bars denote the weaker of the two single-family success rates under the same model and severity; colored bars denote cross-family success; the right-side numbers show the mean difference between cross-family and weakest-component success.}
    \label{fig:cross-vs-single}
\end{figure}

A3+C4 is the strongest negative interaction. Under medium severity, its average cross-family success rate is 0.259, below the weakest single-family component at 0.545; under heavy severity, it drops to 0.027, below the weakest component at 0.250. This combination couples inter-tool confusability with return-evidence competition, so errors no longer stay within a single stage: boundary instability in tool selection and corrupted evidence in observation recovery jointly surface as incorrect arguments. Its failure composition reflects this mechanism, with 88\% incorrect-argument failures.

B4+C4 is the second strong negative interaction. Its cross-family success is 0.357 under medium severity and 0.054 under heavy severity, compared with weakest-component baselines of 0.545 and 0.250. User-request conflict already requires the model to identify the primary target, while evidence competition adds conflicting returned evidence. The combined setting therefore amplifies intent ambiguity into observation-driven incorrect arguments: 73\% of failures are incorrect arguments, and 20\% are tool misselection.

Not every composition creates additional loss. A2+B3 has a heavy cross-family average of 0.304, above the weakest single-family component at 0.190, and B2+D2 has a medium cross-family average of 0.455, above the weakest component at 0.378. We interpret these apparent reversals as floor effects or salience effects rather than true robustness gains: when one single-family component is already near the floor, additional perturbations may not produce much measurable loss, and some runtime failures can become more explicit under combined conditions.

The main conclusion is that cross-family difficulty is not reducible to the weakest single-family component. Tool-output-centered combinations are especially damaging because they convert earlier interface or intent uncertainty into corrupted evidence selection. Full success matrices, failure attribution, and model diagnostic collages are reported in Appendix Figures~\ref{fig:app-cross-success}, \ref{fig:app-cross-failure}, \ref{fig:app-cross-model-sheet}, and \ref{fig:app-gpt-cross-matrix}.

\section{Discussion}

Two takeaways follow from the results. First, clean tool-calling ability cannot replace robustness evaluation. Multiple models are close to perfect under clean conditions, yet their perturbed macro averages and family profiles differ substantially. A single clean score therefore hides whether a model is fragile in schema grounding, argument binding, evidence recovery, or runtime-state judgment.

Second, observation and runtime feedback are current bottlenecks. Tool-output / observation perturbation is the largest-drop family for every model, and runtime-environment perturbation exposes failures in deciding status, retryability, and final result. Agent systems should therefore treat tool returns as structured evidence with provenance, freshness, and machine-decidable status fields, rather than handing unverified logs or natural-language summaries directly to the model.

Third, cascade-aware attribution changes how failures should be used for model and system improvement. If only the observed label is recorded, a downstream incorrect-argument symptom may lead developers to tune argument extraction, even when the primary source is tool-boundary confusion or schema grounding. By separating primary errors from observed symptoms, \ToolRobustBench{} makes failure analysis more actionable: interface designers can improve tool naming and descriptions when S1/S2 failures dominate, while runtime designers can harden status schemas and retry semantics when S4 failures dominate.

Fourth, cross-family evaluation should be treated as a necessary complement to single-family diagnosis. Single-family perturbations identify where a model is fragile in isolation, but deployed agents often face simultaneous interface ambiguity, indirect user requests, noisy returned evidence, and unstable execution feedback. The strong negative interactions in A3+C4 and B4+C4 show that robustness risks can be amplified when earlier-stage uncertainty meets corrupted evidence. A practical evaluation suite should therefore include both stage-isolated tests for localization and a small set of mixed perturbations for interaction stress testing.

\section{Conclusion}

\ToolRobustBench{} provides a stage-aligned, deterministic, and cascade-aware benchmark for tool-calling robustness. Across 15,456 single-family instances and representative mixed perturbations, clean success is near saturated but robustness drops sharply, especially under tool-output / observation perturbation. Cross-family results show that multi-source difficulty cannot be inferred from the weakest single-family component alone. The benchmark therefore supports not only model ranking, but also diagnosis of where failures begin, how they cascade, and which tool-interface or runtime designs are likely to improve robust agent behavior.

\section{Limitations}

The benchmark is intentionally controlled. The main single-family results cover 14 subtypes, but the main text reports only five representative cross-family combinations rather than an exhaustive ranking of all pairwise templates. The evaluated tools are deterministic local tools, which improve reproducibility but do not cover all live API ecosystems. Tasks are single-step and do not yet test long-horizon planning, persistent state, or cross-tool dependency. Model calls use different providers or OpenAI-compatible backends, so protocol compatibility remains a residual source of variation. Finally, the paper treats raw-record recomputation as the standard and documents scorer normalization, compatibility checks, and data-freezing issues in the appendix.

\bibliography{custom}

\clearpage
\appendix
\setcounter{figure}{6}

\section{Tool Environment}

\begin{table*}[htbp]
\centering
\small
\begin{tabularx}{\linewidth}{l r Y}
\toprule
Tool group & \# tools & Tool IDs \\
\midrule
\texttt{calendar\_ops} & 4 & \texttt{cancel\_calendar\_event}, \texttt{check\_calendar\_availability}, \texttt{create\_calendar\_event}, \texttt{update\_calendar\_event} \\
\texttt{commerce\_ops} & 4 & \texttt{compare\_product\_price}, \texttt{estimate\_shipping}, \texttt{lookup\_product}, \texttt{track\_order} \\
\texttt{conversion} & 2 & \texttt{convert\_currency}, \texttt{convert\_unit} \\
\texttt{crm\_ops} & 4 & \texttt{create\_support\_ticket}, \texttt{lookup\_customer}, \texttt{lookup\_invoice}, \texttt{update\_ticket\_status} \\
\texttt{database\_ops} & 4 & \texttt{aggregate\_table}, \texttt{lookup\_record}, \texttt{query\_table}, \texttt{validate\_record} \\
\texttt{email\_ops} & 4 & \texttt{classify\_email}, \texttt{draft\_email}, \texttt{search\_email}, \texttt{summarize\_email} \\
\texttt{file\_ops} & 4 & \texttt{convert\_file\_format}, \texttt{extract\_file\_metadata}, \texttt{search\_file}, \texttt{summarize\_file} \\
\texttt{list\_ops} & 2 & \texttt{filter\_items}, \texttt{sort\_items} \\
\texttt{route\_travel} & 2 & \texttt{calculate\_route}, \texttt{estimate\_travel\_time} \\
\texttt{search\_retrieval} & 3 & \texttt{search\_document}, \texttt{search\_faq}, \texttt{search\_policy} \\
\texttt{security\_ops} & 4 & \texttt{audit\_login\_events}, \texttt{check\_permission}, \texttt{revoke\_session}, \texttt{rotate\_api\_key} \\
\texttt{weather\_env} & 3 & \texttt{get\_air\_quality}, \texttt{get\_temperature}, \texttt{get\_weather} \\
\bottomrule
\end{tabularx}
\caption{Complete deterministic local-tool environment in \ToolRobustBench. The main experiments sample 16 tools from this 40-tool environment.}
\label{tab:tool-env}
\end{table*}

\section{Evaluation Validity Details}
\label{app:validity}

The main findings are robust to scorer normalization and protocol-compatibility adjustments. For runtime-environment perturbations, failure-type alias normalization prevents semantically equivalent labels from being marked wrong only because of surface-form mismatch. After normalization, the D-family scores remain differentiated: \texttt{deepseek-v4-pro} reaches 0.748, \texttt{claude-sonnet-4-6} reaches 0.745, \texttt{qwen-plus} reaches 0.713, and \texttt{gpt-5.4-mini} reaches 0.708. This indicates that normalization fixes a scorer artifact without eliminating model differences.

For post-execution decisions in the tool-output and runtime-environment families, compatibility checks ensure that empty or invalid structured outputs are not caused by backend token-budget or protocol failures. Strict raw records remain the traceable main results; compatibility-adjusted estimates, when used, should be reported separately as sensitivity analysis. Perturbation calibration figures and scorer-ablation plots are provided in Appendix Figures~\ref{fig:app-alias}, \ref{fig:app-nonmonotonic}, and \ref{fig:app-gemini-profile}.

\section{Evidence Table}

\begin{table*}[htbp]
\centering
\small
\begin{adjustbox}{width=\linewidth}
\begin{tabular}{p{0.28\linewidth} p{0.20\linewidth} p{0.25\linewidth} p{0.07\linewidth} p{0.12\linewidth}}
\toprule
Data source & Citable metrics & Corresponding conclusion & Main text & Figure \\
\midrule
All-family raw records & sample count, success, failure tag, stage fields & Seven models complete full raw records, totaling 15,456 instances & Yes & Main Figures~\ref{fig:family-heatmap}--\ref{fig:failure-taxonomy}; appendix diagnostics \\
Overall summaries & condition-level summary & Most summaries can be used as references, but Claude D has inconsistency & Cautiously & No \\
Family summaries & family/severity success & Mostly consistent with raw records; main-text values use raw recomputation & Yes & Figure~\ref{fig:family-heatmap}, Appendix Figure~\ref{fig:app-severity} \\
Subtype summaries & subtype success & C2, D2, and B3 are key difficult subtypes and support non-monotonic response analysis & Yes & Appendix Figures~\ref{fig:app-subtype}, \ref{fig:app-nonmonotonic}, \ref{fig:app-gemini-profile} \\
Failure-taxonomy records & primary failure tag & The main failure mode of each family is highly concentrated & Yes & Figure~\ref{fig:failure-taxonomy} \\
Boundary reports & cascade rate, boundary-violation rate & Cascade failure and boundary violation need separate explanations & Yes & Appendix Figure~\ref{fig:app-cascade-boundary} \\
Cross-family raw records & cross-family success, failure tag, cascade fields & Medium/heavy interaction results for five representative pairwise combinations & Yes & Figure~\ref{fig:cross-vs-single}, Appendix Figures~\ref{fig:app-cross-success}--\ref{fig:app-cross-failure} \\
Cross-vs-single summary & cross success, weakest single, delta vs. weakest & Comparison between cross-family and corresponding single-family components & Yes & Figure~\ref{fig:cross-vs-single} \\
Cross-family plot manifest & retained plot type, source plot directory & Consistency of retained plot types for each model $\times$ combination & No & Appendix Figures~\ref{fig:app-cross-model-sheet}--\ref{fig:app-gpt-cross-matrix} \\
Evaluation-validity notes & scorer correction and compatibility diagnosis & Evaluation-validity analysis and ablation design & Yes & Appendix Figure~\ref{fig:app-alias} \\
Run manifests & sampled tools, dataset scope, family config & 16 sampled tools, complete environment of 40 tools/12 groups, single-step setting & Yes & Appendix Table~\ref{tab:tool-env} \\
Figure assets & main figures & The main text keeps four core figures and moves secondary figures to the appendix & Yes & Main and appendix figures \\
\bottomrule
\end{tabular}
\end{adjustbox}
\caption{Evidence table linking data sources, citable metrics, conclusions, and figures.}
\label{tab:evidence}
\end{table*}

\section{Complete Model Result Tables}

Table~\ref{tab:family-overall} summarizes the perturbed macro-average results of each model over the four families.

\begin{table*}[htbp]
\centering
\small
\begin{tabular}{lrrrrr}
\toprule
Model & Tool Interface & User Intent & Tool Output & Runtime Environment & Overall \\
\midrule
\texttt{claude-sonnet-4-6} & 0.968 & 0.785 & 0.568 & 0.745 & 0.766 \\
\texttt{deepseek-v4-pro} & 0.926 & 0.809 & 0.512 & 0.748 & 0.749 \\
\texttt{gpt-5.4-mini} & 0.856 & 0.781 & 0.503 & 0.708 & 0.712 \\
\texttt{qwen-plus} & 0.931 & 0.802 & 0.380 & 0.713 & 0.706 \\
\texttt{gemini-2.5-pro} & 0.900 & 0.748 & 0.481 & 0.637 & 0.692 \\
\texttt{deepseek-v4-flash} & 0.958 & 0.783 & 0.319 & 0.616 & 0.669 \\
\texttt{gemini-2.5-flash} & 0.889 & 0.700 & 0.420 & 0.648 & 0.664 \\
\bottomrule
\end{tabular}
\caption{Family-level robustness for each model.}
\label{tab:family-overall}
\end{table*}

Tables~\ref{tab:claude-severity}--\ref{tab:qwen-severity} give the success rate of each model under family $\times$ severity. All values are recomputed directly from raw \texttt{records.csv}.

\begin{table}[htbp]
\centering
\small
\begin{tabular}{lrrrr}
\toprule
Family & Clean & Light & Medium & Heavy \\
\midrule
Interface & 1.000 & 1.000 & 0.972 & 0.931 \\
Intent & 1.000 & 0.974 & 0.870 & 0.510 \\
Output & 1.000 & 0.661 & 0.625 & 0.417 \\
Runtime & 1.000 & 0.972 & 0.708 & 0.556 \\
\bottomrule
\end{tabular}
\caption{Family $\times$ severity success rates for \texttt{claude-sonnet-4-6}.}
\label{tab:claude-severity}
\end{table}

\begin{table}[htbp]
\centering
\small
\begin{tabular}{lrrrr}
\toprule
Family & Clean & Light & Medium & Heavy \\
\midrule
Interface & 1.000 & 1.000 & 0.972 & 0.903 \\
Intent & 0.979 & 0.953 & 0.839 & 0.557 \\
Output & 0.750 & 0.406 & 0.370 & 0.182 \\
Runtime & 0.979 & 1.000 & 0.521 & 0.326 \\
\bottomrule
\end{tabular}
\caption{Family $\times$ severity success rates for \texttt{deepseek-v4-flash}.}
\label{tab:deepseek-flash-severity}
\end{table}

\begin{table}[htbp]
\centering
\small
\begin{tabular}{lrrrr}
\toprule
Family & Clean & Light & Medium & Heavy \\
\midrule
Interface & 1.000 & 1.000 & 0.924 & 0.854 \\
Intent & 1.000 & 0.990 & 0.885 & 0.552 \\
Output & 1.000 & 0.719 & 0.589 & 0.229 \\
Runtime & 1.000 & 1.000 & 0.729 & 0.514 \\
\bottomrule
\end{tabular}
\caption{Family $\times$ severity success rates for \texttt{deepseek-v4-pro}.}
\label{tab:deepseek-pro-severity}
\end{table}

\begin{table}[htbp]
\centering
\small
\begin{tabular}{lrrrr}
\toprule
Family & Clean & Light & Medium & Heavy \\
\midrule
Interface & 0.958 & 0.986 & 0.875 & 0.806 \\
Intent & 0.958 & 0.901 & 0.823 & 0.375 \\
Output & 0.979 & 0.635 & 0.495 & 0.130 \\
Runtime & 0.958 & 0.958 & 0.431 & 0.556 \\
\bottomrule
\end{tabular}
\caption{Family $\times$ severity success rates for \texttt{gemini-2.5-flash}.}
\label{tab:gemini-flash-severity}
\end{table}

\begin{table}[htbp]
\centering
\small
\begin{tabular}{lrrrr}
\toprule
Family & Clean & Light & Medium & Heavy \\
\midrule
Interface & 1.000 & 0.979 & 0.903 & 0.819 \\
Intent & 1.000 & 0.932 & 0.875 & 0.438 \\
Output & 0.979 & 0.693 & 0.562 & 0.188 \\
Runtime & 0.938 & 0.958 & 0.486 & 0.465 \\
\bottomrule
\end{tabular}
\caption{Family $\times$ severity success rates for \texttt{gemini-2.5-pro}.}
\label{tab:gemini-pro-severity}
\end{table}

\begin{table}[htbp]
\centering
\small
\begin{tabular}{lrrrr}
\toprule
Family & Clean & Light & Medium & Heavy \\
\midrule
Interface & 1.000 & 1.000 & 0.840 & 0.729 \\
Intent & 0.979 & 0.969 & 0.875 & 0.500 \\
Output & 0.979 & 0.703 & 0.641 & 0.167 \\
Runtime & 0.979 & 0.938 & 0.618 & 0.569 \\
\bottomrule
\end{tabular}
\caption{Family $\times$ severity success rates for \texttt{gpt-5.4-mini}.}
\label{tab:gpt-mini-severity}
\end{table}

\begin{table}[htbp]
\centering
\small
\begin{tabular}{lrrrr}
\toprule
Family & Clean & Light & Medium & Heavy \\
\midrule
Interface & 1.000 & 1.000 & 0.944 & 0.847 \\
Intent & 1.000 & 0.984 & 0.896 & 0.526 \\
Output & 1.000 & 0.526 & 0.417 & 0.198 \\
Runtime & 1.000 & 0.854 & 0.729 & 0.556 \\
\bottomrule
\end{tabular}
\caption{Family $\times$ severity success rates for \texttt{qwen-plus}.}
\label{tab:qwen-severity}
\end{table}

\begin{table*}[htbp]
\centering
\small
\begin{tabular}{lrrrrr}
\toprule
Family & Failures & Cascade rate & Expected-primary match & Spillover match & Boundary violation \\
\midrule
Tool Interface & 247 & 0.352 & 0.891 & 0.352 & 0.109 \\
User Intent & 917 & 0.108 & 0.957 & 0.108 & 0.043 \\
Tool Output & 2198 & 0.003 & 0.931 & 0.000 & 0.069 \\
Runtime Environment & 944 & 0.638 & 0.994 & 0.638 & 0.006 \\
\bottomrule
\end{tabular}
\caption{Family-level cascade and boundary diagnostic statistics.}
\label{tab:boundary-stats}
\end{table*}

\FloatBarrier
\section{Appendix Figure Index}

For review and Overleaf upload, main-text figures are stored in \texttt{picture/} and ordered from \texttt{fig01\_} to \texttt{fig06\_}. The main text retains four core figures: Figure~\ref{fig:pipeline}, which explains the evaluation pipeline; Figure~\ref{fig:family-heatmap}, which shows family-level differences; Figure~\ref{fig:failure-taxonomy}, which supports failure attribution; and Figure~\ref{fig:cross-vs-single}, which supports the core comparison between cross-family and single-family baselines. The original leaderboard and cascade-boundary figures are retained as appendix figures.

Appendix figures are stored in \texttt{picture/} and \texttt{picture/appendix/}. The detailed figures moved down from the original main text cover dataset coverage, clean-vs-robustness, severity curves, subtype difficulty, post-execution diagnostics, alias ablation, non-monotonic cases, Gemini-Pro execution profile, cross-family success matrix, and cross-family failure diagnostics.

Figures \texttt{app\_fig13}--\texttt{app\_fig18} are appendix collages and complete cross-family diagnostic figures. They supplement family-level diagnostic sheets, cross-family model overviews, and the full pairwise composition space for \texttt{gpt-5.4-mini}.

\ifincludeappendixfigures
\begin{figure*}[t]
\centering
\safeincludegraphics[width=0.95\linewidth]{picture/fig02_overall_leaderboard.png}
\caption{Overall robustness leaderboard with family decomposition.}
\label{fig:app-leaderboard}
\end{figure*}

\begin{figure*}[t]
\centering
\safeincludegraphics[width=0.95\linewidth]{picture/appendix/app_fig02_dataset_coverage.png}
\caption{Dataset coverage.}
\label{fig:app-dataset-coverage}
\end{figure*}

\begin{figure*}[t]
\centering
\safeincludegraphics[width=0.95\linewidth]{picture/appendix/app_fig03_clean_vs_robustness.png}
\caption{Clean capability versus robustness.}
\label{fig:app-clean-robustness}
\end{figure*}

\begin{figure*}[t]
\centering
\safeincludegraphics[width=0.95\linewidth]{picture/appendix/app_fig04_severity_degradation.png}
\caption{Severity response curves.}
\label{fig:app-severity}
\end{figure*}

\begin{figure*}[t]
\centering
\safeincludegraphics[width=0.95\linewidth]{picture/appendix/app_fig05_subtype_difficulty.png}
\caption{Subtype difficulty.}
\label{fig:app-subtype}
\end{figure*}

\begin{figure*}[t]
\centering
\safeincludegraphics[width=0.95\linewidth]{picture/appendix/app_fig06_execution_stage2_diagnostics.png}
\caption{Post-execution diagnostics.}
\label{fig:app-stage2}
\end{figure*}

\begin{figure*}[t]
\centering
\safeincludegraphics[width=0.95\linewidth]{picture/fig05_cascade_boundary_diagnostics.png}
\caption{Cascade and boundary diagnostics by perturbation family.}
\label{fig:app-cascade-boundary}
\end{figure*}

\begin{figure*}[t]
\centering
\safeincludegraphics[width=0.95\linewidth]{picture/appendix/app_fig08_alias_normalization_ablation.png}
\caption{Failure-type alias normalization ablation.}
\label{fig:app-alias}
\end{figure*}

\begin{figure*}[t]
\centering
\safeincludegraphics[width=0.95\linewidth]{picture/appendix/app_fig09_nonmonotonic_severity_cases.png}
\caption{Non-monotonic severity-response cases.}
\label{fig:app-nonmonotonic}
\end{figure*}

\begin{figure*}[t]
\centering
\safeincludegraphics[width=0.95\linewidth]{picture/appendix/app_fig10_gemini_pro_execution_profile.png}
\caption{Gemini-Pro execution profile.}
\label{fig:app-gemini-profile}
\end{figure*}

\begin{figure*}[t]
\centering
\safeincludegraphics[width=0.95\linewidth]{picture/appendix/app_fig11_crossfamily_selected_success_matrix.png}
\caption{Cross-family success matrix.}
\label{fig:app-cross-success}
\end{figure*}

\begin{figure*}[t]
\centering
\safeincludegraphics[width=0.95\linewidth]{picture/appendix/app_fig12_crossfamily_failure_diagnostics.png}
\caption{Cross-family failure diagnostics.}
\label{fig:app-cross-failure}
\end{figure*}

\begin{figure*}[t]
\centering
\safeincludegraphics[width=0.95\linewidth]{picture/appendix/app_fig13_interface_exposure_diagnostic_sheet.png}
\caption{Tool-interface diagnostic sheet.}
\label{fig:app-interface-sheet}
\end{figure*}

\begin{figure*}[t]
\centering
\safeincludegraphics[width=0.95\linewidth]{picture/appendix/app_fig14_user_intent_diagnostic_sheet.png}
\caption{User-intent diagnostic sheet.}
\label{fig:app-intent-sheet}
\end{figure*}

\begin{figure*}[t]
\centering
\safeincludegraphics[width=0.95\linewidth]{picture/appendix/app_fig15_observation_diagnostic_sheet.png}
\caption{Tool-output / observation diagnostic sheet.}
\label{fig:app-observation-sheet}
\end{figure*}

\begin{figure*}[t]
\centering
\safeincludegraphics[width=0.95\linewidth]{picture/appendix/app_fig16_execution_environment_diagnostic_sheet.png}
\caption{Runtime-environment diagnostic sheet.}
\label{fig:app-execution-sheet}
\end{figure*}

\begin{figure*}[t]
\centering
\safeincludegraphics[width=0.95\linewidth]{picture/appendix/app_fig17_crossfamily_model_diagnostic_sheet.png}
\caption{Cross-family model diagnostic sheet.}
\label{fig:app-cross-model-sheet}
\end{figure*}

\begin{figure*}[t]
\centering
\safeincludegraphics[width=0.95\linewidth]{picture/appendix/app_fig18_gpt54mini_full_crossfamily_composition_matrix.png}
\caption{Full cross-family composition matrix for \texttt{gpt-5.4-mini}.}
\label{fig:app-gpt-cross-matrix}
\end{figure*}
\fi

\end{document}
